% GENERATED FILE. Edit canonical lessons, then run npm run generate:books.
% canonical-source-hash: fnv1a64:d00d715585b293f8
% canonical-lessons: KA-C121-topi, KA-C121-ungura, KA-C121-bale, KA-C121-saru, KA-C121-sambaru

\chapter{Cap, Ring, Bangle, Rasam, Sambar}
\label{ch:ka-cap-ring-bangle-rasam-sambar}
\hlchaptermodality{121}

\begin{chapteropening}
\textbf{By the end of this chapter:} \emph{I can say a cap, a ring, a bangle, rasam and sambar.}

Complete the last lesson of chapter 121: I can say a cap, a ring, a bangle, rasam and sambar.
\end{chapteropening}

\section[\texorpdfstring{\kn{ಟೋಪಿ} (ṭōpi)}{ಟೋಪಿ (ṭōpi)}]{\kn{ಟೋಪಿ} (ṭōpi) — a cap, a hat}

\label{lesson:KA-C121-topi}



\begin{quote}
Before the new one: say the Kannada for a dhoti, then the Kannada for sandals.
\end{quote}

\subsection*{You'll want to know: \kn{ಟೋಪಿ}}
\textbf{\kn{ಟೋಪಿ}} — \emph{ṭōpi} — \textquotedblleft{}a cap, a hat\textquotedblright{}.

\begin{cousinweb}
\textbf{\kn{ಟೋಪಿ}} (\emph{ṭōpi}) came from Hindustani \emph{ṭōpī}, \textquotedblleft{}cap\textquotedblright{}. Its deeper origin is not settled.
\end{cousinweb}

\begin{grammarlens}[title={What you've built}]
One more word for this chapter.
\end{grammarlens}

\subsection*{Your turn}
\begin{itemize}
\raggedright
  \item \emph{Say it:} \emph{ṭōpi}
  \item \emph{Say it:} \emph{ṭōpi}, once more
  \item \emph{Say it:} say \emph{ṭōpi}
  \item \emph{Recall:} say the Kannada for a dhoti, then the Kannada for sandals, then say \emph{ṭōpi} again
\end{itemize}

\begin{tcolorbox}[breakable,colback=teal!4,colframe=teal!35!black,title={Before you move on}]
What does \kn{ಟೋಪಿ} mean? (\textquotedblleft{}A cap, a hat\textquotedblright{}.) Say it once more.
\end{tcolorbox}
\section[\texorpdfstring{\kn{ಉಂಗುರ} (uṅgura)}{ಉಂಗುರ (uṅgura)}]{\kn{ಉಂಗುರ} (uṅgura) — a ring}

\label{lesson:KA-C121-ungura}



\begin{quote}
Before the new one: say the Kannada for sandals, then the Kannada for a cap.
\end{quote}

\subsection*{You'll want to know: \kn{ಉಂಗುರ}}
\textbf{\kn{ಉಂಗುರ}} — \emph{uṅgura} — \textquotedblleft{}a ring\textquotedblright{}.

\begin{cousinweb}
\textbf{\kn{ಉಂಗುರ}} (\emph{uṅgura}) is native Dravidian, one word with Telugu \emph{uṅgaramu} (DEDR 572).
\end{cousinweb}

\begin{grammarlens}[title={What you've built}]
One more word for this chapter.
\end{grammarlens}

\subsection*{Your turn}
\begin{itemize}
\raggedright
  \item \emph{Say it:} \emph{uṅgura}
  \item \emph{Say it:} \emph{uṅgura}, once more
  \item \emph{Say it:} say \emph{uṅgura}
  \item \emph{Recall:} say the Kannada for sandals, then the Kannada for a cap, then say \emph{uṅgura} again
\end{itemize}

\begin{tcolorbox}[breakable,colback=teal!4,colframe=teal!35!black,title={Before you move on}]
What does \kn{ಉಂಗುರ} mean? (\textquotedblleft{}A ring\textquotedblright{}.) Say it once more.
\end{tcolorbox}
\section[\texorpdfstring{\kn{ಬಳೆ} (baḷe)}{ಬಳೆ (baḷe)}]{\kn{ಬಳೆ} (baḷe) — a bangle}

\label{lesson:KA-C121-bale}



\begin{quote}
Before the new one: say the Kannada for a cap, then the Kannada for a ring.
\end{quote}

\subsection*{You'll want to know: \kn{ಬಳೆ}}
\textbf{\kn{ಬಳೆ}} — \emph{baḷe} — \textquotedblleft{}a bangle\textquotedblright{}.

\begin{cousinweb}
\textbf{\kn{ಬಳೆ}} (\emph{baḷe}) is native Dravidian, one word with Tamil \emph{vaḷai} and Telugu \emph{balayu} (DEDR 5313). DEDR lists \emph{baḷe}, \textquotedblleft{}ring, armlet, bracelet\textquotedblright{}, under the root for \textquotedblleft{}to go round\textquotedblright{}. DEDR 5313 compares Sanskrit \emph{val}.
\end{cousinweb}

\begin{grammarlens}[title={What you've built}]
One more word for this chapter.
\end{grammarlens}

\subsection*{Your turn}
\begin{itemize}
\raggedright
  \item \emph{Say it:} \emph{baḷe}
  \item \emph{Say it:} \emph{baḷe}, once more
  \item \emph{Say it:} say \emph{baḷe}
  \item \emph{Recall:} say the Kannada for a cap, then the Kannada for a ring, then say \emph{baḷe} again
\end{itemize}

\begin{tcolorbox}[breakable,colback=teal!4,colframe=teal!35!black,title={Before you move on}]
What does \kn{ಬಳೆ} mean? (\textquotedblleft{}A bangle\textquotedblright{}.) Say it once more.
\end{tcolorbox}
\section[\texorpdfstring{\kn{ಸಾರು} (sāru)}{ಸಾರು (sāru)}]{\kn{ಸಾರು} (sāru) — rasam, a thin curry}

\label{lesson:KA-C121-saru}



\begin{quote}
Before the new one: say the Kannada for a ring, then the Kannada for a bangle.
\end{quote}

\subsection*{You'll want to know: \kn{ಸಾರು}}
\textbf{\kn{ಸಾರು}} — \emph{sāru} — \textquotedblleft{}rasam, a thin curry\textquotedblright{}.

\begin{cousinweb}
\textbf{\kn{ಸಾರು}} (\emph{sāru}) is native Dravidian, one word with Tamil \emph{cāṟu} (DEDR 2484). DEDR lists Kannada \emph{sāṟu}, \textquotedblleft{}juice, broth, pepper water\textquotedblright{}.
\end{cousinweb}

\begin{grammarlens}[title={What you've built}]
One more word for this chapter.
\end{grammarlens}

\subsection*{Your turn}
\begin{itemize}
\raggedright
  \item \emph{Say it:} \emph{sāru}
  \item \emph{Say it:} \emph{sāru}, once more
  \item \emph{Say it:} say \emph{sāru}
  \item \emph{Recall:} say the Kannada for a ring, then the Kannada for a bangle, then say \emph{sāru} again
\end{itemize}

\begin{tcolorbox}[breakable,colback=teal!4,colframe=teal!35!black,title={Before you move on}]
What does \kn{ಸಾರು} mean? (\textquotedblleft{}Rasam, a thin curry\textquotedblright{}.) Say it once more.
\end{tcolorbox}
\section[\texorpdfstring{\kn{ಸಾಂಬಾರು} (sāmbāru)}{ಸಾಂಬಾರು (sāmbāru)}]{\kn{ಸಾಂಬಾರು} (sāmbāru) — sambar}

\label{lesson:KA-C121-sambaru}



\begin{quote}
Before the new one: say the Kannada for a bangle, then the Kannada for rasam.
\end{quote}

\subsection*{You'll want to know: \kn{ಸಾಂಬಾರು}}
\textbf{\kn{ಸಾಂಬಾರು}} — \emph{sāmbāru} — \textquotedblleft{}sambar\textquotedblright{}.

\begin{cousinweb}
The origin of \textbf{\kn{ಸಾಂಬಾರು}} (\emph{sāmbāru}) is uncertain. Neither DEDR nor CDIAL records it, and no source for it is settled.
\end{cousinweb}

\begin{grammarlens}[title={What you've built}]
One more word for this chapter.
\end{grammarlens}

\subsection*{Your turn}
\begin{itemize}
\raggedright
  \item \emph{Say it:} \emph{sāmbāru}
  \item \emph{Say it:} \emph{sāmbāru}, once more
  \item \emph{Say it:} say \emph{sāmbāru}
  \item \emph{Recall:} say the Kannada for a bangle, then the Kannada for rasam, then say \emph{sāmbāru} again
\end{itemize}

\begin{tcolorbox}[breakable,colback=teal!4,colframe=teal!35!black,title={Before you move on}]
What does \kn{ಸಾಂಬಾರು} mean? (\textquotedblleft{}Sambar\textquotedblright{}.) Say it once more.
\end{tcolorbox}
